\chapter{Introduction}
\label{chapter:introduction}

\section{Background and Motivation}

Rolling-element bearings support shafts, carry loads, and reduce friction in motors, pumps, turbines, gearboxes, fans, compressors, and many other machines. Because they sit directly in the load path, a bearing defect rarely remains an isolated issue for long. A small defect on a raceway or rolling element can create repeated impacts, excite vibration, produce heat, damage nearby parts, and eventually force the machine out of service. The consequences of a bearing failure are sometimes far greater than the cost of the bearing in industrial maintenance.

While traditional maintenance approaches could help, they don't completely address the issue. Reactive maintenance is waiting for failure and takes the risk of unplanned shutdown. Preventive maintenance involves exchanges of parts at a predetermined interval, typically a more conservative interval than necessary. Can take out healthy components too soon or miss faults that occur quicker than they should. The improvements of both methods are provided in condition-based maintenance, in that the state of the machine is assessed and a maintenance decision is made based on the evidence. Prognostics and health management is a further step that attempts to predict the remaining useful life of a component before failure, in addition to predicting failure, \cite{jardine2006,lei2018}.

The reason why the remaining useful life prediction for bearings is difficult is that the deterioration process is not linear. Vibration may stay low for a considerable length of the test, but then increase dramatically towards the end of life. The measured signals are affected by the type of fault, the speed of the fault, the load, the position of the sensor, noise, resonance, lubrication, and the geometry of the bearings \cite{wu2022, kannan2024}. Training on one of the bearings doesn't mean that the model works well on another bearing, even if they are both from the same test rig. This is the reason that the IMS bearing data set is helpful for benchmarking run-to-failure experiments. They enable a study of the same raw measurements and the comparison of methods under identical conditions \cite{qiu2006}.

The deep learning techniques have gained popularity for their ability to capture nonlinear patterns within vibration signals, which is crucial for bearing diagnosis and RUL prediction. CNNs can learn local patterns from windows of data, LSTMs can model temporal dependence, and feed-forward networks can fit engineered feature vectors \cite{babu2016,berghout2022,ayman2025}. These methods can work well when the training and test data are similar. Their weakness is that they usually learn only from examples. When labeled run-to-failure data are limited or the test bearing degrades differently from the training bearings, a data-only model may produce predictions that are numerically plausible but physically questionable.

Physics-informed neural networks offer a way to include prior engineering knowledge in the learning process. In their most established form, PINNs train a neural network while penalizing violations of governing equations or physical constraints \cite{raissi2019}. Bearing RUL prediction does not usually provide a complete closed-form degradation equation, but it does provide useful partial knowledge. Damage should generally increase with operating time, RUL should generally decrease, and bearing geometry gives characteristic fault frequencies associated with cage, inner-race, outer-race, and rolling-element defects. This thesis studies whether such weak physical information can improve RUL prediction on the IMS data.

\section{Problem Statement}

The problem studied in this thesis is the prediction of normalized remaining useful life for selected bearings in the IMS run-to-failure dataset. The main question is whether a physics-informed neural network using time and fault-frequency information can generalize better across bearing runs than comparable data-driven neural baselines.

This question is practical as well as methodological. In a real machine, the future failure path is unknown. A maintenance engineer may have historical data from one bearing, but the next bearing may not degrade in the same way. A useful RUL model must therefore do more than memorize one run. It must learn a relationship between measured condition and remaining life that transfers, at least partly, across bearings.

\section{Objectives of the Study}

The objectives of this thesis are:

\begin{enumerate}[label=\arabic*.]
    \item To extract degradation-relevant vibration features from selected IMS bearing runs and construct normalized RUL targets with a PCA-based health indicator for analysis.
    \item To design and evaluate a DeepXDE physics-informed neural model that uses weak monotonicity and fault-frequency energy priors for bearing RUL prediction.
    \item To compare the proposed physics-informed model with data-only, LSTM, and CNN neural baselines under complete held-out bearing splits and interpret where each model succeeds or fails.
\end{enumerate}

\section{Scope}

This study is limited to the IMS bearing data placed locally under \texttt{data/raw/1st\_test}, \texttt{data/raw/2nd\_test}, and \texttt{data/raw/3rd\_test}. Four bearing runs are used in the final experiment design: \texttt{ds2\_b1}, \texttt{ds1\_b3}, \texttt{ds1\_b4}, and \texttt{ds3\_b3}. The work uses vibration signals only. It does not add temperature, oil debris, acoustic emission, or motor-current data.

The models predict normalized RUL from extracted features. The thesis does not claim field deployment readiness. The data come from a controlled test rig, and the operational conditions are much simpler than many industrial machines. There is also an intentional weak physical information from PINN. It relies on monotonicity and energy priors for fault frequency, instead of a full energy balance contact mechanics or crack growth model.

\section{Importance of the Study}

This work is the integration of three parts of the bearing prognostics problem in one local workflow: Signal processing, machine learning and physical interpretation. The feature extraction step maintains the links with vibration analysis of RMS, kurtosis, crest factor, and envelope spectral energies. The neural models investigate if those features can contribute to Predicting RULs over bearing runs. The physics-informed model is used to test if simple Priorities can be regularized in the learning process through physical priors.
The final result isn't just a success story for the proposed model. That is important. A thesis result should state what happens during the experiment, not what is expected to happen in the proposal. The final run indicates that LSTM sequence modeling was able to capture the degradation .The proposed weakly physics-informed feed-forward model is less consistent in their trend. Meanwhile, the PINN was competitive in Experiment 2 and ranked first in Experiment 2 (RMSE).This mixed behavior is more clearly specified for future work than an overconfident
positive claim would.

\section{Limitations}

The main limitations are:

\begin{enumerate}[label=\arabic*.]
    \item The IMS dataset has small number of total run-to-failure records which creates a experiment design constraint.
    \item The physics selected prior is not a full bearing degradation law. It is a regularizer based on monotonic RUL behavior and fault-frequency energy.
    \item The model uses hand-engineered features, not raw vibration waveform.
    \item The RUL labels are derived from elapsed time to the end of each run, which assumes the last available file represents failure or near-failure.
    \item The experiments are run on one local PC, so training time and numerical results may vary on other machines.
\end{enumerate}

\section{Thesis Outline}

Chapter~\ref{chapter:literature} reviews bearing failure, vibration-based condition monitoring, RUL prediction, deep learning baselines, and physics-informed learning. Chapter~\ref{chapter:methodology} describes the IMS dataset, feature extraction, preprocessing, health indicator construction, model architectures, and evaluation metrics. Chapter~\ref{chapter:implementation} describes the local implementation, experiment splits, software environment, and reproducibility setup. Chapter~\ref{chapter:results} presents the results and discusses model behavior across the three experiments. Chapter~\ref{chapter:conclusion} summarizes the findings, contributions, limitations, and future work.
